\documentclass[11pt]{article}
\usepackage[margin=1in]{geometry}
\usepackage{graphicx}
\usepackage{float}
\usepackage[hidelinks]{hyperref}
\input{generated/metrics.tex}
\title{probe}
\author{Computational Study Team}
\date{}
\begin{document}
\maketitle
\begin{abstract}
This synthetic validation fixture exercises a reproducible computational research
workflow using a known analytical reference. A fixed-seed Monte Carlo estimate of
$\pi$ from \SampleCount{} points is \PiEstimate{}, with absolute error
\AbsoluteError{}. This is an execution fixture, not a novel research finding.
\end{abstract}
\section{Question and methods}
We sample independent coordinate pairs uniformly in the unit square using Python
3.11.13's seeded pseudorandom generator, with seed \SeedValue{}. Four times the
fraction satisfying $x^2+y^2\leq1$ estimates $\pi$. The reported plug-in standard
error is $4\sqrt{\hat p(1-\hat p)/n}$. This numerical-integration example follows
the general Monte Carlo approach described by Metropolis and Ulam
\cite{metropolis1949}. The sample count and seed are fixed before execution.
\section{Results}
\begin{center}
\input{generated/table.tex}
\end{center}
\begin{figure}[H]
\centering
\includegraphics[width=\linewidth]{figures/convergence.pdf}
\caption{Cumulative estimates at fixed checkpoints. The reference line is the
analytical value of $\pi$; a single trajectory does not establish convergence rates.}
\end{figure}
\section{Reproduction and limitations}
Run \texttt{make reproduce} in a clean checkout to regenerate the experiment,
figures, tables, and this paper. The repository pins the Python interpreter and
Tectonic engine; the run record identifies inputs and output hashes. Reproduction
compares numeric results at prespecified tolerances, not PDF bytes.
This single synthetic run validates the execution pathway only. Its standard
error describes sampling variability, and no significance or novelty claim is
made. An independent evidence review is required before verified study status.
\bibliographystyle{plain}
\bibliography{references}
\end{document}
